\origpage{11}{}
\textbf{4.1} \textit{Let $\Sigma$ be a surface, and let $\Sigma'$ be obtained by cutting along $F \subseteq \Sigma$.}

\hspace*{12pt}1. \textit{If $F$ is a proper line with ends in different cuffs then $\hat{\Sigma}'$ is homeomorphic to $\hat{\Sigma}$ and $c(\delm{\hat{\Sigma}'}\fix{\Sigma'}) = c(\delm{\hat{\Sigma}}\fix{\Sigma}) - 1$.}\fixnote{fix:4.1}{Here $\hat{\Sigma}$ has no boundary, so $c(\hat{\Sigma}) = 0$. What is meant is that $\Sigma'$ has one cuff fewer than $\Sigma$.}

\hspace*{12pt}2. \textit{If $F$ is a proper line with ends in the same cuff and $\Sigma'$ is connected then $\hat{\Sigma}'$ is simpler than $\hat{\Sigma}$.}

\hspace*{12pt}3. \textit{If $bd(\Sigma) = \emptyset$ and $F$ is an $O$-arc and $\Sigma'$ is connected then $\hat{\Sigma}'$ is simpler than $\Sigma$.}

The proof is straightforward\del{, and we omit it}. \fix{We give the details. We use the classification of surfaces: every surface is homeomorphic to the surface obtained from a sphere by adding $a$ handles and $b$ crosscaps and then deleting the interiors of $c$ disjoint closed discs, for some $a, b, c \ge 0$, and then it has $c$ cuffs. In each part, $\Sigma$ is obtained from $\Sigma'$ by pasting together the two copies of $F$ in $bd(\Sigma')$.}

\fix{Let $F$ be a proper line, let $X, X' \subseteq bd(\Sigma')$ be its two copies, which are disjoint lines, and suppose that $\Sigma'$ is connected and is obtained from a sphere by adding $a$ handles and $b$ crosscaps and deleting $c$ discs. If $X$ and $X'$ lie in the same cuff of $\Sigma'$, then pasting $X$ onto $X'$ either splits this cuff into two cuffs of $\Sigma$, each containing one end of $F$, and then $\Sigma$ is obtained in the same way with $a, b, c + 1$; or it leaves one cuff, containing both ends of $F$, and then $\Sigma$ is obtained with $a, b + 1, c$. If $X$ and $X'$ lie in different cuffs of $\Sigma'$, then the pasting joins these two cuffs into one cuff of $\Sigma$, containing both ends of $F$, and $\Sigma$ is obtained with $a + 1, b, c - 1$ or with $a, b + 2, c - 1$.}

\fix{For part 1, let the ends of $F$ be in the cuffs $C_1 \ne C_2$ of $\Sigma$. Then $X$, $X'$ and what is left of $C_1, C_2$ form one cuff of $\Sigma'$; and $\Sigma'$ is connected, since every component of $\Sigma'$ meets $X \cup X'$ (because $\Sigma$ is connected). Hence we are in the first case above, and so $\hat{\Sigma}'$ is homeomorphic to $\hat{\Sigma}$ and $c(\Sigma') = c(\Sigma) - 1$. For part 2, the ends of $F$ are in the same cuff of $\Sigma$, so the first case does not occur, and $\hat{\Sigma}$ is obtained from $\hat{\Sigma}'$ by adding a crosscap, a handle, or two crosscaps. Hence $\hat{\Sigma}'$ is simpler than $\hat{\Sigma}$. For part 3, let $\Sigma'$ be connected and obtained with $a, b, c$. If $F$ has a neighbourhood in $\Sigma$ that is an annulus, the two copies of $F$ are two cuffs of $\Sigma'$, and pasting them together gives a surface with $a + 1$ handles and $b$ crosscaps, or with $a$ handles and $b + 2$ crosscaps. Otherwise $F$ has a neighbourhood that is a M\"obius band, the copies of $F$ form one cuff of $\Sigma'$, and the pasting adds a crosscap. In both cases $\Sigma = \hat{\Sigma}$ is obtained from $\hat{\Sigma}'$ by adding handles or crosscaps, and so $\hat{\Sigma}'$ is simpler than $\Sigma$.}\fixnote{fix:4.1-proof-added}{The paper omits the proof of (4.1). The argument above is the one used in the formalization.}

We shall also have to deal with line and $O$-arcs $F$ which separate $\Sigma$ (that is, such that $\Sigma \setminus F$ is disconnected), but in these cases there are surfaces $\Sigma_1, \Sigma_2 \subseteq \Sigma$ with $\Sigma_1 \cup \Sigma_2 = \Sigma$ and $\Sigma_1 \cap \Sigma_2 = F$, and we can get away with using these subsurfaces. In these cases, therefore, no cutting is needed, which is convenient for purposes of notation. The corresponding results are:

\textbf{4.2} \textit{Let $\Sigma$ be a surface, and let $\Sigma_1, \Sigma_2 \subseteq \Sigma$ be surfaces with $\Sigma_1 \cup \Sigma_2 = \Sigma$ and $\Sigma_1 \cap \Sigma_2 = F$.}

\hspace*{12pt}1. \textit{If $F$ is a proper line in $\Sigma$ with both ends in the same cuff, then either $\hat{\Sigma}_2$ is simpler than $\hat{\Sigma}$ or $\hat{\Sigma}_1$ is a sphere.}

\hspace*{12pt}2. \textit{If $F$ is a proper line in $\Sigma$ with both ends in the same cuff, and $\hat{\Sigma}_1$ is a sphere, then $\hat{\Sigma}_2$ is homeomorphic to $\hat{\Sigma}$; and either $c(\Sigma_2) < c(\Sigma)$, or $c(\Sigma_2) = c(\Sigma)$ and there is a closed disc $\Delta \subseteq \Sigma$ with $F \subseteq bd(\Delta) \subseteq F \cup bd(\Sigma)$.}

\hspace*{12pt}3. \textit{If $bd(\Sigma) = \emptyset$ and $F$ is an $O$-arc in $\Sigma$, then either $\hat{\Sigma}_2$ is simpler than $\hat{\Sigma}$ or $\Sigma_1$ is a closed disc.}

\fix{We give the proofs, with the classification of surfaces as in the proof of (4.1). Let $\Sigma_i$ be obtained from a sphere by adding $a_i$ handles and $b_i$ crosscaps and deleting $c_i$ discs ($i = 1, 2$). In parts 1 and 2, $F$ is a proper line in $\Sigma$ with ends in a cuff $C$ of $\Sigma$; then $F$ lies in one cuff of each of $\Sigma_1, \Sigma_2$, made of $F$ and one of the two lines of $C$ between the ends of $F$, and $\Sigma$ is obtained by pasting $\Sigma_1$ and $\Sigma_2$ together along $F$. This joins these two cuffs into the cuff $C$ and keeps the other cuffs, and so $\Sigma$ is obtained from a sphere by adding $a_1 + a_2$ handles and $b_1 + b_2$ crosscaps and deleting $c_1 + c_2 - 1$ discs. For part 1, if $a_1 = b_1 = 0$ then $\hat{\Sigma}_1$ is a sphere, and otherwise $\hat{\Sigma}$ is obtained from $\hat{\Sigma}_2$ by adding $a_1$ handles and $b_1$ crosscaps, and so $\hat{\Sigma}_2$ is simpler than $\hat{\Sigma}$. For part 2, $a_1 = b_1 = 0$, so $\hat{\Sigma}_2$ is homeomorphic to $\hat{\Sigma}$, and $c(\Sigma_2) = c(\Sigma) - c_1 + 1 \le c(\Sigma)$, since $c_1 \ge 1$. If equality holds then $c_1 = 1$, so $\Sigma_1$ is a closed disc; its boundary is its one cuff, which contains $F$ and is contained in $F \cup bd(\Sigma)$, and we take $\Delta = \Sigma_1$. For part 3, $bd(\Sigma) = \emptyset$, so $F$ is the only cuff of $\Sigma_1$ and of $\Sigma_2$, and pasting them along $F$ shows that $\Sigma$ is obtained from a sphere by adding $a_1 + a_2$ handles and $b_1 + b_2$ crosscaps. If $a_1 = b_1 = 0$ then $\Sigma_1$ is a closed disc, and otherwise $\hat{\Sigma}_2$ is simpler than $\hat{\Sigma}$, as in part 1.}\fixnote{fix:4.2-proof-added}{The paper states (4.2) without proof. The argument above is the one used in the formalization.}

A drawing in a surface $\Sigma$ is a pair $(U, V)$ where $U \subseteq \Sigma$ is closed, $V \subseteq U$ is finite, $U \cap bd(\Sigma) \subseteq V$, $U \setminus V$ has only finitely many connected components, called \textit{edges}, and for each edge $e$, either

\hspace*{12pt}1. its closure $\bar{e}$ is an $O$-arc with $|\bar{e} \cap V| = 1$, or

\hspace*{12pt}2. its closure is a line meeting $V$ in precisely its ends.

If $\Gamma = (U, V)$ is a drawing we write $U(\Gamma) = U$ and $V(\Gamma) = V$. A drawing $\Gamma$ is therefore a graph with vertex set $V(\Gamma)$, and we use graph-theoretic terminology for drawings without further explanation.

If $\Gamma$ is a drawing in $\Sigma$, and $F \subseteq \Sigma$, we say that $F$ is $\Gamma$-\textit{normal} if $F \cap U(\Gamma) \subseteq V(\Gamma)$. If $F$ is a $\Gamma$-normal $O$-arc or line which does not separate $\Sigma$, and $\Sigma'$ is obtained by cutting along $\Sigma$, then
